A task studied in developmental linguistics \cite{carey1978acquiring} is the ability to learn and utilize new words, for example using a word in a sentence after seeing it defined only once, or conversely inferring a word’s meaning from only one usage.  Here we qualitatively test GPT-3’s ability to do the former.  Specifically, we give GPT-3 the definition of a nonexistent word, such as ``Gigamuru'', and then ask it to use it in a sentence.  We provide one to five previous examples of a (separate) nonexistent word being defined and used in a sentence, so the task is few-shot in terms of previous examples of the broad task and one-shot in terms of the specific word.  Table \ref{completion:novel words} shows the 6 examples we generated; all definitions were human-generated, and the first answer was human-generated as conditioning while the subsequent answers were generated by GPT-3.  These examples were generated continuously in one sitting and we did not omit or repeatedly try any prompts.  In all cases the generated sentence appears to be a correct or at least plausible use of the word.  In the final sentence the model generates a plausible conjugation for the word ``screeg'' (namely ``screeghed''), although the use of the word is slightly awkward (``screeghed at each other'') despite being plausible in the sense that it could describe a toy sword fight.  Overall, GPT-3 appears to be at least proficient at the task of using novel words in a sentence.

